\section{Überblick}

Der Standard-Workflow für die forensische Image-Analyse folgt immer demselben Schema:
Integrität prüfen → Image mounten → Partitionen analysieren → Dateien extrahieren → Aufräumen (LIFO).

\subsection{1. Integrität prüfen}

Vor jeder Analyse müssen die Prüfsummen des Images verifiziert werden.
Der gespeicherte MD5-Hash muss mit dem berechneten übereinstimmen.

\begin{lstlisting}[language=bash]
ewfverify image.E01
md5sum image.E01 image.E02 ...
\end{lstlisting}

\subsection{2. Image mounten}

EWF-Image als virtuelles RAW-Device mit Overlay bereitstellen.
Der Parameter \texttt{--cache} ist essenziell: alle Änderungen werden im Overlay gespeichert, das Original bleibt unverändert.

\begin{lstlisting}[language=bash]
mkdir -p /tmp/ewf
xmount --in ewf image.E0* --cache /tmp/img.ovl --out raw /tmp/ewf
\end{lstlisting}

\subsection{3. Partitionstabelle analysieren}

\begin{lstlisting}[language=bash]
mmls /tmp/ewf/image.dd
\end{lstlisting}

Die Ausgabe zeigt das Partitionslayout.
Der \textbf{Start}-Offset jeder Partition wird als \texttt{-o} Parameter für alle TSK-Tools benötigt.

\subsection{4. Loop Device erstellen}

\begin{lstlisting}[language=bash]
losetup --partscan --find --show /tmp/ewf/image.dd
\end{lstlisting}

Erstellt \texttt{/dev/loopXpN} für jede Partition automatisch.
Alternative für eine einzelne Partition:

\begin{lstlisting}[language=bash]
losetup -o $((OFFSET*512)) /dev/loop0 /tmp/ewf/image.dd
\end{lstlisting}

\subsection{5. Dateisystem analysieren}

\begin{lstlisting}[language=bash]
fsstat /dev/loopXpN | head
fls -pr /dev/loopXpN
fls -r -d /dev/loopXpN    # nur geloeschte Dateien
\end{lstlisting}

\begin{itemize}
    \item \texttt{fsstat} — zeigt Dateisystemtyp, Version, Blockgröße
    \item \texttt{fls -pr} — listet alle Dateien rekursiv mit vollem Pfad
    \item Gelöschte Dateien sind mit \texttt{*} markiert
\end{itemize}

\subsection{6. Dateien extrahieren}

\begin{lstlisting}[language=bash]
fls -pr /dev/loopXpN | grep -i 'filename'
icat /dev/loopXpN INODE > /tmp/recovered_file
file /tmp/recovered_file
\end{lstlisting}

Nach jeder Extraktion \textbf{immer} mit \texttt{file} verifizieren.

\subsection{7. Partition mounten (optional)}

\begin{lstlisting}[language=bash]
mount -o ro /dev/loopXpN /mnt
# oder mit Offset:
mount -o ro,loop,offset=$((OFFSET * 512)) /tmp/ewf/image.dd /mnt
\end{lstlisting}

Read-Only-Mount zum Browsen.
Zeigt nur Allocated Space (logische Analyse).

\subsection{8. Aufräumen (LIFO)}

Umgekehrte Reihenfolge — zuletzt gemountet wird zuerst unmountet:

\begin{lstlisting}[language=bash]
umount /mnt
losetup -D
umount /tmp/ewf
\end{lstlisting}
